\documentclass[11pt]{article}
\usepackage[T1]{fontenc}
\usepackage{lmodern,amsmath,amssymb,amsthm,booktabs}
\usepackage[margin=1in]{geometry}
\usepackage[hidelinks]{hyperref}
\newtheorem{proposition}{Proposition}
\title{Rectangle incidence circuits for a stronger integer-multiplication exponent}
\author{Douglas Colkitt\thanks{Prepared with assistance from OpenAI Codex.
This is not independent review or endorsement by OpenAI.}}
\date{October 7, 2026}
\begin{document}
\maketitle
\begin{abstract}
We replace individual neighbor channels in the finite bit network of OpenAI's
integer-multiplication manuscript by rectangle circuits with shared intermediate
computations. Nondegenerate common frames preserve the absolute rank deficit.
Together with nonadjacent axis routing and the previously audited parameter
recipe, the construction supports the conditional saving $\kappa=2^{-67}$,
a factor $256$ above our preceding $2^{-75}$ witness. The complex network is
unchanged. No practical runtime improvement or global optimality is claimed.
\end{abstract}

\section{Scope and construction}
The model is a fixed finite-alphabet Turing machine with a fixed finite number
of one-dimensional tapes. We assume the unchanged algorithmic interfaces of
\emph{Integer multiplication below $n\log n$}, pinned at
\begin{center}\texttt{adc7f1241b42e322a6451854ab7e4b4c146bf78a}.\end{center}
The construction below replaces the bit primitive. The remaining parts of
the manuscript's multiplication theorem are assumed.
\subsection{A rectangle circuit for the bit network}
\label{sec:incidence-circuit}
Retain $h=46$, $\mathcal T=\binom{[h]}3$, $v=15180$, $N=v^3$, and
$m=h^3=97336$. We change only the bit network. The complex network and its
guard constants are unchanged. On $F=\mathbb Q^h$ use the form $I-J/9$;
the indicator $t_T$ has norm two, and neighboring indicators are orthogonal.

Every ordered neighbor pair has a unique common point $i$. Removing $i$
leaves an ordered pair of disjoint two-subsets of the other 45 points.
Partition this disjointness relation into rectangles $\mathcal S\times
\mathcal U$. In each rectangle every source subset is disjoint from every
target subset. Restoring $i$ gives source and target triple families, all
cross pairs of which are neighbors.

Here is an explicit finite procedure for the partition used below. Let
$D_{a,b}(n)$ be the relation between disjoint $a$- and $b$-subsets, with
$0\le a,b\le2$. For a partition let $(C,L,R)$ count its rectangles,
source memberships, and target memberships. An empty relation contributes
$(0,0,0)$; if $a=0$ or $b=0$, one full rectangle contributes
$(1,\binom na,\binom nb)$. Available nontrivial choices are singleton-row
stars, singleton-column stars, or any split $n=l+(n-l)$, $1\le l<n$.
For a split, separate pairs by their source and target sizes in the first
part. Each case is a product
$D_{x,y}(l)\times D_{a-x,b-y}(n-l)$.
Product partitions multiply the three counts coordinatewise; disjoint cases
add them. Thus the procedure always produces an exact partition.

Select a plan by minimizing $\alpha L+\beta R-C$, initially
$\alpha=\beta=1$. A product with a complete factor having counts $(1,L_0,R_0)$
passes weights $(\alpha L_0,\beta R_0)$ to the other factor. If both factors
are nontrivial, both are $D_{1,1}$: try for each its selected unweighted plan,
row stars, and column stars, and take the least product cost. Try all split
sizes and the two star alternatives; ties retain the first plan in the order
just specified, with row stars before column stars and increasing split sizes.
This is a finite recursive construction, not a global optimality claim.
The companion integer implementation and exhaustive pair check give
\[
 (C,L,R)=(2759,27278,27534),\qquad L+R-C=52053.
\]
Its rectangles contain each of the
$\binom{45}2\binom{43}2=893970$ ordered disjoint pairs exactly once.

For a rectangle with $a$ sources and $b$ targets allocate $a+b-1$ auxiliary
roles: $a$ input slots and $b$ output slots sharing one pivot. Let $L$ now
denote the invertible bit mixer that adds the other input slots into the
pivot and then adds the pivot into the other output slots. Its output-by-input
submatrix is all ones. Its inverse reverses these two operations.
Let $V$ copy sources into their input slots and $J$ inject output slots into
their targets. Then $JLV$ is exactly this rectangle's neighbor contribution.
Use separate roles for distinct rectangles. The side-role count per invocation is
\[
 R_{\rm side}=46(52053)=2394438.
\]

\subsection{Restoration and common frames}
Let $G,R$ denote the original central gather and scatter. The forward
chronological schedule is
\[
 L,\ J,\ L^{-1},\ R,\ V,\ G,\ R,\ L,\ J,\ L^{-1},\ G,\ V.
\]
Here $L$ is applied independently on each rectangle and all other operations
are their indicated bit additions. For arbitrary initial side vector $z$,
the first three operations add $JLz$ to the target and restore $z$.
The copy changes it to $z+Vx$; the later mixer and injection add
$JL(z+Vx)$. Their dirty-scratch contributions cancel, and the last inverse
mixer and copy restore $z$. The centers contribute $RGx$ and are also restored.
Since $JLV$ is precisely the neighbor matrix,
$(JLV+RG)x=x$. Thus the invocation implements $y\gets y+x$ on arbitrary
inputs. Reversing and inverting its gates implements the inverse invocation.

At stage $j$ write, as in the original tensor construction,
\[
 A=F^{\otimes(j-1)}=B\mathbin\perp P,\quad
 P=\langle t_{a_1}\otimes\cdots\otimes t_{a_{j-1}}\rangle,
 \quad Q=\langle t_{a_{j+1}}\otimes\cdots\otimes t_{a_3}\rangle.
\]
Suppress the future factor $Q$ and abbreviate the following labels:
\begin{align*}
 D_0&=B\otimes F,& D_1&=A\otimes F,\\
 X_t&=(B\otimes F)\mathbin\perp(P\otimes\langle t\rangle),&
 Y_t&=(B\otimes F)\mathbin\perp(P\otimes t^\perp),\\
 D_U&=(B\otimes F)\mathbin\perp(P\otimes U).
\end{align*}
For a forward rectangle $U$ is the span of its physical $X$ indicators.
Every indicator contains its common point $i$. For any vector in this span,
$\sum_j u_j=3u_i$, whence
\[
 \langle u,u\rangle=\sum_{j\ne i}u_j^2>0\quad(u\ne0).
\]
Indeed zero outside coordinates also force $u_i=0$. Thus $U$ is
nondegenerate, including when the ambient form is indefinite. Orthogonality
of all cross pairs gives $\langle t_X\rangle\subset U\subset t_Y^\perp$.

The common frames of the forward gates, in the displayed chronological order,
are
\[
 D_0,\ D_0,\ D_0,\ D_0,\ X_{t_X},\ D_1,\ D_0,\
 D_U,\ Y_{t_Y},\ D_1,\ D_1,\ D_1.
\]
The gates at $X_{t_X}$ and $Y_{t_Y}$ are grouped separately by data triple;
the middle mixers at $D_U$ are grouped by rectangle. The early side injection
is instead at the common label $D_0$, allowing the earlier scratch mixers.
This is a monotone enlargement of the incoming physical $Y$ label.

For stage 2 reverse and invert the scalar schedule with logical source $Y$
and target $X$. Its operations and common frames are respectively
\begin{gather*}
 V,\ G,\ L,\ J,\ L^{-1},\ R,\ G,\ V,\ R,\ L,\ J,\ L^{-1},\\
 D_0,\ D_0,\ D_0,\ X_{t_X},\ D_U,\ D_1,\ D_0,\
 Y_{t_Y},\ D_1,\ D_1,\ D_1,\ D_1.
\end{gather*}
Now $U$ is the span of the physical $X$ indicators, which are the rectangle's
logical targets. The same nondegeneracy and orthogonality proof applies.

All frames above are tensored with $Q$. An input-only side role follows
$D_0\subset X_t\subset D_U\subset D_1$ in a forward invocation;
an output-only role follows $D_0\subset D_U\subset Y_t\subset D_1$;
the pivot follows both. In the reverse invocation the two roles exchange
these physical descriptions. All side edges are therefore increasing, with
nondegenerate orthogonal residuals. Physical $X$ still changes from
$A\otimes\langle t\rangle\otimes Q$ to $(A\otimes F)\otimes Q$,
and physical $Y$ from $B\otimes\langle t\rangle\otimes Q$ to
$Y_t\otimes Q$. These are the original interstage boundaries.
Each central role has one decreasing edge, from $D_1$ to $D_0$, losing
dimension $h$. Thus the total decreasing dimension remains
\[
 L_{\rm dec}=3v^2h^2.
\]

\subsection{Sharing all first-stage and third-stage auxiliary roles}
Pair the ground points consecutively. Define a permutation $\pi$ of triples
as follows. A triple containing a full pair and a singleton keeps the singleton
and cycles its full pair among the pairs other than the singleton's pair.
A triple using three different pairs keeps the point in the least indexed
pair and flips the other two points to their partners. The first rule is a
cycle on each fixed-singleton class and the second an involution on the
remaining classes. Hence $\pi$ is bijective and $|T\cap\pi(T)|=1$.

Match the whole auxiliary bank of the stage-1 invocation with fixed triples
$(A,B)$ to that of stage 3 with fixed triples $(B,\pi(A))$, keeping the local
role index. Both invocations restore arbitrary auxiliary inputs, and the
matching is bijective. No simultaneous roles are identified. It includes
centers as well as side roles.

Every auxiliary first gate now has label $D_0\otimes Q$ and every last gate
has label $D_1\otimes Q$. In a matched pair the joining labels are
\[
 E=F\otimes\langle t_A\rangle\otimes\langle t_B\rangle,
 \qquad
 H=\langle t_B\otimes t_{\pi(A)}\rangle^\perp_{F\otimes F}\otimes F.
\]
Both are nondegenerate and $E\subset H$, because $t_A\perp t_{\pi(A)}$.
The edge $E\to H$ replaces $E\to F^{\otimes3}$ and $0\to H$,
reducing total projection-difference rank by exactly $m$ per removed role.
No additional decreasing edge arises. The final role count and rank sum are
\begin{align*}
 W_{\rm b}^{\rm rect}&=2N+2v^2(R_{\rm side}+h)=1110529317427200,\\
 s_{\rm b}^{\rm rect}&=W_{\rm b}^{\rm rect}m-N+2L_{\rm dec}
                         =108094481068699857600,\\
 W_{\rm b}^{\rm rect}m-s_{\rm b}^{\rm rect}&=572394081600>0.
\end{align*}

\begin{proposition}[Rectangle bit interface]\label{prop:incidence-bit-interface}
The rectangle network exchanges the two scalar bit banks and fixes every
auxiliary input. Its rational frames satisfy the original all-role endpoint
identity, with $W_{\rm b}^{\rm rect}$ roles and total edge rank
$s_{\rm b}^{\rm rect}<W_{\rm b}^{\rm rect}m$.
\end{proposition}
\begin{proof}
The forward, inverse, forward stages give the scalar bank exchange. The
restoration proof permits the auxiliary identifications. Every surviving
auxiliary role still starts with matrix zero and ends with $I_m$ and is fixed
by that exchange; data-role endpoints are unchanged. Thus the original
endpoint identity holds on every role. All side and data changes are nested;
only the stated central decreases remain. Telescoping gives total absolute
label change $W_{\rm b}^{\rm rect}m-2N+2L_{\rm dec}$, and the unchanged
rational source correction adds $N$. The common-frame identity therefore
applies with the displayed rank sum. The mixers, rational bases, and shear
factorizations are all fixed finite data, independent of the eventual input
length. The original fixed-tape interchange interface consequently applies.
\end{proof}

\subsection{Explicit rational bounds for the two recurrences}
Put
\[
 a_{\rm b}=\frac{46}{10^{11}},\qquad
 a_{\rm c}=\frac9{500000000000},\qquad L_0=\frac{5743}{500}.
\]
The new bit deficit and retained complex deficit satisfy
\[
 \eta_{\rm b}^{\rm rect}=\frac{27}{5098849024}>a_{\rm b}L_0,
 \qquad \eta_{\rm c}=\frac7{22253827054}>a_{\rm c}L_0.
\]
For $1\le x\le2$, write $z=(x-1)/(x+1)$ and
\[
 S(x)=2\sum_{j=0}^{23}\frac{z^{2j+1}}{2j+1},\qquad
 E(x)=\frac{2z^{49}}{49(1-z^2)}.
\]
The logarithm series bounds $S(x)\le\log x\le S(x)+E(x)$.
Exact rational comparison gives
$\log m\le16(S(2)+E(2))+S(m/2^{16})+E(m/2^{16})<L_0$.
Therefore $m^{1-a}>m(1-aL_0)>m(1-\eta)$ for either corresponding pair.
We may take
\begin{equation}\label{eq:explicit-motif-exponents}
 \tau=1-a_{\rm b},\qquad\sigma=1-a_{\rm c},
\end{equation}
with $s_{\rm b}^{\rm rect}/W_{\rm b}^{\rm rect}<m^\tau$ and
$s_{\rm c}/W_{\rm c}<m^\sigma$.


\section{Conditional multiplication exponent}
Put $a=a_{\rm b}$ and use the nonadjacent-axis parameter recipe
\begin{gather*}
 \epsilon=1/21,\quad \beta=9/10,\quad c=9a/10,\quad
 \lambda=1-19a^2/20,\quad\lambda'=1-9a^2/10,\\
 \delta=1/16,\qquad C_1=20.
\end{gather*}
The packed comparison is $\tau(1+c/\beta)=1-a^2<\lambda<\lambda'<1$.
The separate complex conditions follow from $19a^2/20<a_{\rm c}$ and
$9a^2<a_{\rm c}$. Exact arithmetic gives
\[
 \min_j g_j=g_2=g_3=\frac{3a^2}{70}
 =\frac{1587}{175000000000000000000000}>2^{-67}.
\]
The other margins are $g_1>1/2$, $g_4=20a/21$, $g_5=13/112$,
$g_6=299/336$, and $g_7=1/21$, all strictly larger. The positive gap absorbs
the fixed logarithmic factors as in the existing strict-margin argument,
yielding $O(n(\log n)^{1-2^{-67}})$ conditionally.

The dimension exponent, stopping threshold, complex network, and guard
analysis are retained: the guard is $p^{20/21}=o(p)$, the new chunk has
$K=\Theta(p^{3a/70})=o(\ell)$ with $\ell=\Theta(p^{20/21})$, and
$K/\log p\to\infty$. The numerical maps and final rounding are unchanged.
No use is made of the separate recurrence refinement requiring
$\sigma\le\tau$; the original separate layer inequalities pass here.

\section{Reproduction and boundary of verification}
The companion artifacts are:
\begin{itemize}
\item \path{scripts/incidence_rectangles.py} and \path{scripts/incidence_network.py};
\item \path{certificates/incidence-network.json};
\item \path{patches/h46-incidence-67.patch}.
\end{itemize}
The rectangle checker enumerates every
ordered disjoint pair at the target size. Small forward and inverse invocations
are checked on every input basis vector, including every scratch input. A full
three-stage small network checks bank exchange and restoration with shared
nonzero scratch. Rational tests verify the norm identity and sharing inclusion.
The written frame argument and the unchanged upstream machinery remain
mathematical proof obligations; these checks are not a formal proof of the
complete algorithm.
\end{document}
